\section{Interoperability}\label{sec:interop}

A GPU-native logic engine is only as useful as its ability to exchange data with the frameworks researchers already use. xlog exposes GPU-resident relations, probabilities, and gradients through four standard interfaces---DLPack, Arrow, Python bindings, and PyTorch autograd---so that results flow into downstream tools without device-to-host copies or format conversion.

\subsection{DLPack: Zero-Copy GPU Tensor Sharing}

The \texttt{xlog-cuda} crate implements the DLPack~\cite{dlpack} protocol in \texttt{crates/xlog-cuda/src/dlpack.rs}, providing bidirectional zero-copy exchange of GPU buffers with any framework that speaks DLPack---including PyTorch, JAX, CuPy, and cuDF. Export is handled by \texttt{DlpackTable}, which wraps a reference-counted \texttt{CudaBuffer} and produces per-column \texttt{DLManagedTensor} handles with a registered deleter callback; the underlying GPU memory remains shared and is freed only when all consumers release their handles. Import is handled by \texttt{CudaKernelProvider::from\_dlpack\_tensors}, which accepts a vector of \texttt{DLManagedTensor} pointers---one per column---validates device affinity, dtype, contiguity, and alignment, and assembles them into a \texttt{CudaBuffer} without any device-to-host copy. A schema-checked variant (\texttt{from\_dlpack\_tensors\_with\_schema}) adds compile-time type verification against an expected \texttt{Schema}, rejecting mismatches before evaluation begins. All xlog scalar types---\texttt{U32}, \texttt{U64}, \texttt{I32}, \texttt{I64}, \texttt{F32}, \texttt{F64}, \texttt{Bool}, and \texttt{Symbol}---have well-defined DLPack dtype mappings. The protocol is surfaced to Python as PyCapsule objects with the \texttt{"dltensor"} name; consumers call \texttt{torch.from\_dlpack(capsule)} to obtain a live CUDA tensor backed by the same device memory that xlog wrote.

\subsection{Arrow IPC: Columnar Export with Dictionary Encoding}

For integration with columnar analytics tools---Pandas, Polars, DuckDB---xlog exports relations through the Arrow~\cite{arrow} C Data Interface. The \texttt{crates/xlog-cuda/src/arrow\_device.rs} module defines \texttt{ArrowDeviceArray}, a \texttt{repr(C)} struct that pairs Arrow FFI array and schema pointers with a CUDA device type and ordinal, following the Arrow C Device specification. The Python API exposes \texttt{export\_arrow\_device} and \texttt{import\_arrow\_device} for zero-copy round-tripping between xlog's GPU buffers and Arrow-compatible consumers. Symbol-typed columns receive special treatment: \texttt{xlog-core}'s \texttt{symbol::to\_arrow} function (\texttt{crates/xlog-core/src/symbol.rs}) converts interned symbol IDs into Arrow \texttt{DictionaryArray<UInt32Type>} with a \texttt{StringArray} dictionary, preserving the human-readable string mapping while keeping the integer representation compact. The inverse path, \texttt{symbol::from\_arrow}, re-interns dictionary values back into xlog's global symbol table. This dictionary encoding means that a Polars or DuckDB consumer can read symbolic relation columns as native string dictionaries without a separate lookup step.

\subsection{Python Bindings}

The \texttt{pyxlog} crate provides Python bindings built with PyO3~\cite{pyo3}. The public API---documented in \texttt{crates/pyxlog/python/pyxlog/\_native.pyi}---exposes three compilation entry points (\texttt{LogicProgram.compile}, \texttt{Program.compile}, \texttt{IlpProgramFactory.compile}), each returning a compiled program object with evaluation, training, and result-extraction methods. Type stubs (\texttt{.pyi} files) provide full IDE auto-completion and static type checking. All tensor-valued attributes---\texttt{EvalResult.prob}, \texttt{McDeviceEvalResult.query\_counts}, \texttt{IlpTaggedCreditDeviceResult} fields---are returned as DLPack capsules, keeping the data on the GPU and letting the caller choose how to materialize it. Long-running GPU operations release the Python GIL via \texttt{py.allow\_threads()} before entering CUDA kernel dispatch---circuit evaluation, ILP fixpoint iteration, and the neural backward pass all execute without holding the GIL, allowing Python data-loader threads and asyncio tasks to run concurrently with GPU computation.

\subsection{PyTorch Autograd Integration}

Neural-symbolic training requires gradients to flow from logical inference back into neural network parameters. xlog achieves this by constructing loss tensors that participate directly in PyTorch's autograd graph. The \texttt{forward\_backward\_tensor} method in \texttt{crates/pyxlog/src/neural.rs} evaluates the compiled circuit on the GPU---with the GIL released---then exports the scalar NLL loss as a DLPack capsule and imports it into PyTorch via \texttt{torch.from\_dlpack}. Because the neural network's forward pass produced grad-enabled output tensors that fed into the circuit, the resulting loss tensor remains connected to the autograd graph. The batched path (\texttt{forward\_backward\_batch\_complex\_tensor}) goes further: it stacks per-network inputs into a single batched forward pass, evaluates multiple queries against the shared circuit, accumulates per-query losses, and then calls \texttt{torch.autograd.backward(outputs, grad\_tensors)} for a single backward pass through the entire computation graph. This design means that xlog's probabilistic circuit acts as a differentiable function inside PyTorch---gradients flow from the logic layer through DLPack tensors back to \texttt{nn.Module} parameters, with no custom autograd \texttt{Function} subclass required.
